\section{Validation}\label{sec:validation}

\paragraph{Principle.} Every implementation is compared with one reference: a
row-at-a-time, two-pass evaluation of \cref{eq:attention} in double precision (S0).
It computes all scores of a query, their maximum, and then the normalised sum; it
uses neither the online update nor the tuned kernel. It does share the tensor
layout, the head mapping and the mask predicate with the code under test, so a
mistake in those would be common to both; the real-data runs add a check from
outside, PyTorch's own output (\cref{sec:real}). Tolerances were fixed before
running: an absolute difference of at most $10^{-5}$ for the synthetic outputs,
whose components are at most one in magnitude (they are weighted averages of values
in $[-1,1)$), and any NaN or infinity is a failure. Every timed scaling run of this
report repeats the comparison on evenly spaced sample rows, each against all $N$
keys (64 rows in the CPU runs and 32 in the GPU runs unless stated), and checks all
outputs for non-finite values; a failed check makes the program exit with an error
and marks the result line. The sampled rows include the first and the last row of
the sequence but are not placed on shard boundaries. Because every sampled row is
compared against \emph{all} keys, a wrong key or value at any shard boundary would
show in each of them; the boundaries on the query side (uneven and empty shards,
partial tiles) are what the unit tests check on every row. The runs that measure something other than the attention
time carry no comparison: peak memory, the hardware-counter runs of
\cref{tab:align}, and the network microbenchmark.

\paragraph{What the tests cover, and why.} The unit tests are few, eight kinds, and
each targets one way the algorithm can fail (\cref{tab:tests}). A test is repeated
for the parameter values on which that failure depends (both kernels, several tile
sizes, rank counts), because an error at a tile or shard boundary shows only when
the boundary falls in the right place; that is the only reason the number of
individual comparisons is larger than the number of tests.

\begin{table}[t]
\centering\small
\begin{tabularx}{\textwidth}{@{}l Y Y@{}}
\toprule
Test & Cases & Failure it would expose\\
\midrule
Block algebra & key-tile sizes 1, 8, 37, 64 with 37 keys, both kernels & wrong rescaling when the running maximum changes between blocks\\
Uneven sizes & $N=1000$ with tiles of 64; $P=3,5,7$; more ranks than rows & remainder handling, empty shards\\
Masks & causal, window, sink+window; blocks that are entirely masked; reversed block order & NaN from $-\infty-(-\infty)$, off-by-one at the diagonal\\
Stressed scores & queries multiplied by 50 (weights span 30 orders of magnitude) & overflow or loss of the normalisation\\
Extreme tails & a value of $10^{38}$ behind a weight of $e^{-100}$ & a tiny weight treated as if it were larger\\
Grouped heads & $H=14$, $H_{kv}=2$ against explicitly replicated key/value heads & wrong head mapping\\
Every MPI variant & seven shapes $\times$ M0--M8 and three fetch orders, three consecutive runs on one object & wrong owner or offset, stale buffers, missing synchronisation\\
Counters & useful pairs and received bytes against closed forms & skipped or double-counted work\\
\bottomrule
\end{tabularx}
\caption{Unit tests of the CPU code (\code{tests/}). All are compared with the
double-precision reference on every row. The analysis program of \cref{sec:real} has
its own test against an independent recomputation and two constructed edge cases
(\code{tests/analyze\_test.py}).}
\label{tab:tests}
\end{table}

\paragraph{Results.} The unit tests were run on Vanda on exactly the submitted
source (\cref{tab:provenance}). The core tests (\n{unit.core.cases} comparisons once
the kernels and tile sizes are multiplied out) and the MPI tests for
$P=1,\dots,\n{unit.mpi.pmax}$ with two threads per rank (\n{unit.mpi.groups}
combinations of shape, variant and fetch order for each $P$) pass; the
largest absolute error over all of them is $\n{unit.maxerr}$ (in the
\n{unit.maxerr.case} case), below the tolerance by a factor of \n{unit.margin}. The
three CUDA kernels pass \n{gpuunit.configs} boundary, mask and grouped-head cases on
the A40 with every row compared (largest error $\n{gpuunit.maxerr}$). The MPI+CUDA
ring was checked at three levels (\cref{sec:gpuring}): two ranks sharing one GPU,
two GPUs in one node, and one GPU in each of two nodes; all repetitions pass. In the
formal runs every timed repetition validates itself, since the time of a wrong
result is worthless: \ValidatedReps{} repetitions carried a validation and
\ValidatedFailed{} exceeded the tolerance. All of these belong to one deliberately
hard study, described next; no repetition of a scaling experiment failed. Running
the same problem with different rank counts and variants changes the result only in
the last bits: in the strong-scaling runs the error of one problem against double
precision differs by at most $\n{val.rankspread}$ between configurations.
Repeating one configuration, by contrast, gave the same result every time. Floating-point
addition is not associative, so a parallel sum can change from run to run when
its order depends on timing; here it does not. A rank folds the key blocks in the
order fixed by its schedule and not in the order of arrival, threads own disjoint
query rows, and no floating-point reduction runs across ranks. Only a change of
the configuration ($P$, the partition, the variant, the fetch order) regroups the
terms. The unit tests
require the output of three consecutive runs to be identical bit for bit for
every variant. In the formal runs, the repetitions of each of the \RepeatedRuns{}
program runs with more than one validated repetition report the same error to the
last recorded digit, and so do the \RestartedConfigurations{} configurations that
were started more than once, \RestartedOtherJob{} of them in different jobs.

\paragraph{How large is the error, and where does it come from?} For $N$ up to
\n{err.rows.max} the tuned kernel was compared with double precision on
\emph{every} row, for key tiles of \n{err.bk.list} keys. With the ordinary synthetic
data the largest difference is $\n{err.plain.max}$ for every size and tile. With
the queries multiplied by \n{err.alpha}, so that the softmax puts nearly all weight
on a few keys, it is $\n{err.stress.N1024}$ at $N=1{,}024$, $\n{err.stress.N4096}$ at
$N=4{,}096$ and $\n{err.stress.N16384}$ at $N=16{,}384$. The last two exceed the
tolerance fixed in advance. I report them as failures (\n{err.failed.configs}
configurations, \n{err.failed.reps} repetitions) and did not widen the tolerance.

The block-wise update does not explain these failures: the error differs by less
than one per cent between tiles of 16 keys, where a query goes through hundreds of
online updates, and tiles of 4{,}096 keys, where for $N=4{,}096$ all keys form a
single block. (That single block still uses the same score loop and the same
exponential, so the comparison isolates the update, not the whole kernel.) What
the numbers are consistent with is the rounding of the single-precision scores. A
score of magnitude $|s|$ carries a rounding error of order $10^{-7}|s|$, and the
exponential turns it into a relative error of the weight of the same size. The
scores of this test have a standard deviation of $50/3\approx17$ by construction,
so the largest of them is near one hundred, which gives $10^{-5}$; the maximum over
more rows then grows slowly with $N$. A check on a development machine points the
same way: keeping the single-precision scores and evaluating everything after them
in double precision left the error unchanged. I did not repeat that ablation on the
cluster, so this is a supported explanation and not a proof that no
single-precision algorithm could meet the tolerance. The conclusion for use is that
the stated tolerance holds for well-scaled inputs and should be read relative to
the score magnitude for sharply peaked ones, which is how the real tensors of
\cref{sec:real} are checked.

\paragraph{A bug that the tolerance alone did not catch.} The first version of the
vectorised exponential (\cref{sec:serial}) clamped arguments below $-87$ to $-87$
instead of letting the result underflow. With order-one values this is invisible,
since a weight of $e^{-87}$ is negligible. A code review constructed finite inputs
for which it is not: with $q=1$, keys $0$ and $-100$ and values $0$ and $10^{38}$,
the correct output is $10^{38}e^{-100}=3.7\times10^{-6}$ and the clamped kernel returned
$10^{38}e^{-87}=1.65$. The
kernel now detects such ranges and evaluates them with the library exponential, and
the case is part of the unit tests. The episode is the reason the test list
contains ``extreme tails'': a test with well-scaled data would never have found it.
